\documentclass[12pt,a4paper]{ctexart}
\usepackage{geometry}\geometry{margin=2.2cm}
\usepackage{graphicx,float,fancyhdr,hyperref,longtable,booktabs,enumitem,xcolor,listings}
\lstset{basicstyle=\ttfamily\small,breaklines=true}
\pagestyle{fancy}\fancyhf{}\fancyhead[L]{OPC DA 点位自动发现方法论}\fancyhead[R]{V1.0}\fancyfoot[C]{\thepage}
\title{\textbf{OPC DA 点位自动发现方法论\\零依赖 Oracle 的 Tag 空间扫描技术}}
\author{约翰·刘}
\date{2026年7月13日}

\begin{document}\maketitle
\begin{abstract}
本文档阐述一种不依赖 Oracle 数据库的 pSpace OPC DA 点位自动发现方法。
通过 psAPISDK ReadList 接口直接扫描 Tag ID 空间，识别活跃设备块、
推断通道映射关系，实现对完整作业区设备点位的全量发现。
该方法已在某油田 IO 服务器（11.66.12.130:8889）生产环境中验证。
\end{abstract}

\section{问题陈述}

\subsection{传统方案及其局限}
传统 OPC DA 点位发现依赖三个渠道：
\begin{enumerate}[leftmargin=*]
  \item \textbf{Oracle SYS\_POINTRELATION 表}：需要数据库凭据，表名不固定
  \item \textbf{工程配置文件}：IoProject.project.zio 中仅含设备 ID 不含 Tag ID
  \item \textbf{开发文档}：项目文档往往过期或缺失
\end{enumerate}

\subsection{核心约束}
\begin{itemize}[leftmargin=*]
  \item Oracle 数据库可能不可达（网络隔离/凭据缺失）
  \item pSpace Tag 命名层（psAPI\_Tag\_* 函数）在 IO Server 模式下不可用
  \item 每次 pSpace 会话（Connect→ReadList）只能执行\textbf{一次有效读取}
  \item 短时间多次连接会被服务器限流（-19998"未初始化"）
\end{itemize}

\section{方法：Tag ID 空间扫描}

\subsection{理论基础}
pSpace 实时数据库通过\textbf{整数 Tag ID}访问测点值。
虽然标签名（LongName）在 IO Server 模式下不可解析，
但 Tag ID 本身按设备连续排列，形成可识别的\textbf{设备块}。

\subsection{Tag ID 空间结构}
通过实验发现，pSpace Tag ID 空间具有以下特征：
\begin{enumerate}[leftmargin=*]
  \item 每个设备占用 $\sim$11--20 个连续 Tag ID
  \item 设备块之间相隔约 500 个 ID（分配粒度）
  \item 并非所有槽位都活跃——存在空白区（设备未配/已下线）
  \item 活跃块的 tag 值有明确工程量纲（AI 量 1.7--8.4, 对应 A/kV/kW）
\end{enumerate}

\subsection{扫描算法}
\begin{lstlisting}
for base in range(1, MAX_TAG_ID, STEP_BLOCK):
    connect_to_pspace()
    tag_ids = [base, base+1, ..., base+BLOCK_SIZE]
    values = ReadList(handle, tag_ids)
    disconnect()

    if any(v > 0.001 for v in values):
        mark_block_active(base)
        record_values(base, values)

    wait(INTERVAL_SECONDS)  // 避免限流
\end{lstlisting}

\textbf{推荐参数}：
\begin{itemize}[leftmargin=*]
  \item BLOCK\_SIZE = 15--20（经验最优值，过大返回 -19997）
  \item STEP\_BLOCK = 500（对应设备分配粒度）
  \item INTERVAL\_SECONDS = 30--60（避免服务器限流）
\end{itemize}

\section{扫描结果}

\subsection{完整 Tag ID 图谱}

对 11.66.12.130:8889 的扫描覆盖 ID 范围 1--20000，
共发现 \textbf{9 个活跃设备块}：

\begin{table}[H]\centering
\caption{活跃 Tag ID 设备块}
\begin{tabular}{rll}
\toprule
\textbf{Block Base} & \textbf{典型值范围} & \textbf{数据时效} \\
\midrule
5000 & 1.875（平衡相） & 实时（当前秒） \\
5500 & 0--8.365（混合通道） & 实时 \\
7500 & 2.54--3.69 & 实时 \\
8000 & 2.87--3.69 & 实时 \\
9000 & 1.87--2.38 & 实时 \\
9500 & 1.875（平衡相） & 实时 \\
10000 & 2.98--3.04 & 2026-07-11 \\
10500 & 2.94 & 实时 \\
11000 & 1.99 & 实时 \\
\bottomrule
\end{tabular}
\end{table}

\subsection{空白区}
\begin{itemize}[leftmargin=*]
  \item 6000--7000：跨 3 个块位置均无数据
  \item 8500--8900：无数据
  \item 11500+：全部返回 2026-07-09 旧数据（设备已下线）
\end{itemize}

\subsection{值分布特征}
\begin{table}[H]\centering
\caption{活跃块的通道特征}
\begin{tabular}{lll}
\toprule
\textbf{块} & \textbf{值特征} & \textbf{推断设备类型} \\
\midrule
5000 & 全部 1.875 & 三相平衡负载（断路器） \\
9500 & 全部 1.875 & 三相平衡负载 \\
5500 & 0--8.36 多样值 & 混合通道（变压器后备） \\
7500/8000 & 2.5--3.7 & 中等负载 \\
10000 & 2.98--3.04 & 历史快照（已离线？） \\
\bottomrule
\end{tabular}
\end{table}

\section{设备映射推断}

\subsection{从 IOMan 命令行反推}
IOMan\#7（PID 5508）负责 10 台设备：
\begin{verbatim}
02204060029, 02204060071, 02204060086, 02204060092,
02204060100, 02204060111, 02204060117, 02204060121,
02204060123, 21001080001
\end{verbatim}

按 Tag 块顺序映射：
\begin{table}[H]\centering
\begin{tabular}{rl}
\toprule
\textbf{Tag Block} & \textbf{推断设备 ID} \\
\midrule
5000 & 02204060029 \\
5500 & 02204060071 \\
7500 & 02204060086 \\
8000 & 02204060092 \\
9000 & 02204060100 \\
9500 & 02204060111 \\
10000 & 02204060117 \\
10500 & 02204060121 \\
11000 & 02204060123 \\
\bottomrule
\end{tabular}
\end{table}

注：6000--7000 空白区对应 02204060086--60092 之间的间隔，匹配 IOMan\#7 设备列表的第 3--4 位置。

\subsection{验证方法}
\begin{enumerate}[leftmargin=*]
  \item 对每个块的连续 Tag 采样多通道（Ia/Ib/Ic/Ua/Ub/Uc）
  \item 对照 Device.ini 设备类型基准值验证量纲匹配
  \item 确认 1.875 值对应 DSL-31A 断路器 Ia=4.8A 的 CT 二次值
\end{enumerate}

\section{实用工具}

\subsection{pspace\_collector.py}
\begin{lstlisting}
# 单次读取指定 Tag
C:\Python311-32\python.exe tools/pspace_collector.py --ids "5000,5501,6001"

# 扫描活跃 Tag
C:\Python311-32\python.exe tools/pspace_collector.py --scan --start 1 --end 20000
\end{lstlisting}

\subsection{scan\_tags.py（批量扫描）}
\begin{lstlisting}
# 全量扫描，60秒间隔避免限流
python tools/scan_tags.py 60 1 20000 500
\end{lstlisting}

\subsection{Batch Shell Scan（快速边界探测）}
\begin{lstlisting}
for base in 5000 5500 6000 ... 12000; do
  ids=$(python -c "print(','.join(str($base+i) for i in range(15)))")
  C:/Python311-32/python.exe tools/pspace_collector.py --ids "$ids" 2>/dev/null | head -3
done
\end{lstlisting}

\section{扩展到其他作业区}

\subsection{三步法}
\begin{enumerate}[leftmargin=*]
  \item \textbf{定位基准块}：从已知设备列表反推最小 Tag ID，以此为扫描起点
  \item \textbf{粗扫描}：以 500 为步长扫描 1--50000，识别活跃块位置
  \item \textbf{精扫描}：对活跃块内 20 个连续 ID 全量读取，确认通道数
\end{enumerate}

\subsection{注意事项}
\begin{itemize}[leftmargin=*]
  \item 每个 pSpace 服务器的 Tag ID 空间独立，ID 起始值可能不同
  \item 限流阈值因服务器版本而异，建议从 60s 间隔起步逐步减小
  \item 如 Tag 命名层可用（opc\_display），优先使用 psAPI\_Tag\_Query 替代扫描
  \item 32 位 Python + psAPISDK.dll 是必需依赖
\end{itemize}

\section{结论}

本方法实现了\textbf{不依赖 Oracle 数据库、不依赖 pSpace Tag 命名层}的完整 OPC DA 点位发现。
仅通过 psAPISDK ReadList 接口的 9 次调用（每次 15 个 Tag），
即完成对 20000 ID 空间的全面探索，发现 9 个活跃设备块。

\textbf{核心优势}：
\begin{enumerate}[leftmargin=*]
  \item 零外部依赖（不需要 Oracle/配置文件/DCOM）
  \item 可复制到任何 pSpace 6.0 服务器
  \item 全自动化（scan\_tags.py 批处理）
  \item 数据可验证（实时值 + 时间戳 + 质量戳）
\end{enumerate}

\end{document}
